\section{Background}
\label{sec:background}

\subsection{Transmission, reporting and two estimation targets}
\label{sec:bg-model}

Consider transmission over one generation, indexed by $t$, in which $n$ infected individuals act as primary cases. Conditional on the instantaneous reproduction number $R_t$, the number of secondary infections caused by primary case $i$, $X_{i,t}$, is assumed to follow a negative binomial distribution, which captures individual heterogeneity in transmission \citep{lloydsmith2005superspreading}. The offspring distribution is given by:
\begin{equation}
  X_{i,t}\mid R_t \sim \NB\!\left(k,\ \frac{k}{k+R_t}\right),\qquad
  \E(X_{i,t}\mid R_t)=R_t,\qquad \Var(X_{i,t}\mid R_t)=R_t+\frac{R_t{}^2}{k},
  \label{eq:offspring}
\end{equation}
where $k>0$ is the dispersion parameter, smaller values of $k$ correspond to greater heterogeneity, and $k\to\infty$ recovers Poisson offspring. Each infection is then reported independently with probability $\rho\in(0,1]$. Because a sum of independent negative binomial variables with a common success probability is negative binomial, and because binomial thinning preserves this family, the reported count $C_t$ satisfies:
\begin{equation}
  C_t\mid n,R_t \sim \NB\!\left(nk,\ \frac{k}{k+\rho R_t}\right),\qquad
  \E(C_t\mid n,R_t)=\rho nR_t,\qquad
  \Var(C_t\mid n,R_t)=\rho nR_t+\frac{\rho^2nR_t{}^2}{k},
  \label{eq:observed}
\end{equation}
where $m=\rho n$ is the expected number of reported cases per period when $R_t=1$. We refer to $m$ as the case volume.

The model rests on four assumptions. Each infection is reported independently with a probability $\rho$ that does not change over time (H1). Primary cases transmit independently and identically according to Eq.~(\ref{eq:offspring}) (H2). The dispersion $k$ does not depend on $R_t$ or $n$ (H3). The number of primary cases $n$ is fixed by design and is independent of $R_t$ (H4). Assumption H4 holds when $n$ is set by contact tracing; in routine surveillance $n$ is itself the outcome of earlier transmission, and the renewal formulation in Section~\ref{sec:meth-loss} relaxes it.

Transmission potential varies between periods with weather, school terms, holidays and behaviour \citep{dalziel2018urbanization}. Its marginal moments are denoted $\E(R_t)=\Rbar$ and $\Var(R_t)=v_R$. Because this variation acts on every chain of transmission in the same period, it corresponds to environmental stochasticity in population ecology, whereas the offspring and reporting variation of Eq.~(\ref{eq:observed}) corresponds to demographic stochasticity \citep{engen1998demographic, lande2003stochastic}.

The scaled count $\Rhat_t=C_t/m$ can be used for two different targets. The first is the period-specific value $R_t$, which is the quantity of interest for situational awareness and for judging whether transmission is currently above one. The second is the mean across periods, $\Rbar$, which is the quantity of interest when describing typical transmission in a season. When $R_t$ is independent across periods, $\Rbar$ is also the best one-step-ahead predictor of next period's value. The error for the first target is $\Rhat_t-R_t$. The error for the second target can be written as:
\begin{equation}
  \Rhat_t-\Rbar=(\Rhat_t-R_t)+(R_t-\Rbar),
  \label{eq:error-split}
\end{equation}
where the second term does not depend on $m$. Conditional on $R_t$, $\Rhat_t$ is unbiased for $R_t$, and its variance $R_t/m+\rho R_t{}^2/(km)$ falls to zero as $m$ grows; in this case $\Rhat_t$ attains the Cram\'er--Rao bound for $R_t$ (Appendix~\ref{app:proofs}).

\subsection{Error variance and the balance scale}
\label{sec:bg-variance}

For the second target, the error variance follows from the law of total variance applied to Eq.~(\ref{eq:error-split}).

\begin{proposition}[Error variance for the mean across periods]
\label{prop:average}
Under H1--H4, the variance of $\Rhat_t=C_t/(\rho n)$ about $\Rbar$ is given by:
\begin{equation}
  V(n,\rho)=v_R+\frac{\Rbar}{\rho n}+\frac{\Rbar^2+v_R}{kn}=v_R+\frac{A}{m},
  \qquad A=\Rbar+\frac{\rho(\Rbar^2+v_R)}{k},
  \label{eq:var-average}
\end{equation}
where the three terms are the environmental, reporting and offspring contributions. Then:
\begin{enumerate}[label=(\roman*),leftmargin=2em]
  \item for fixed $\rho$, $V\to v_R$ as $n\to\infty$;
  \item for fixed $n$, $V\to v_R+\Rbar/n+(\Rbar^2+v_R)/(kn)>v_R$ as $\rho\to1$;
  \item at the balance scale $\mx=A/v_R$, the sampling and environmental contributions are equal, and doubling $m$ reduces $V$ by the fraction $1/[2(1+m/\mx)]$;
  \item among estimators that are conditionally unbiased for $R_t$, no estimator has a smaller mean squared error for $\Rbar$ than $V$.
\end{enumerate}
\end{proposition}

The balance scale plays the role of the critical population size in ecological models, at which demographic and environmental contributions to the variance of population growth are equal \citep{lande2003stochastic}. It is a reference point for comparing the two contributions and does not by itself define a minimum operating size. Part (iv) is restricted to conditionally unbiased estimators. An estimator that shrinks $\Rhat_t$ towards a prior mean can have smaller error for $\Rbar$, and averaging over $T$ independent periods reduces the environmental term to $v_R/T$ (Proposition~\ref{prop:allocation}).

Parts (i) and (ii) show that case volume can be raised in two ways with different limits. Coverage $n$ removes both sampling terms in the limit, whereas full reporting leaves the offspring term. Near the baseline, a small proportional increase in $n$ reduces $V$ by $A/\Rbar$ times as much as the same proportional increase in $\rho$. The limit $v_R$ also assumes that environmental variation is independent between periods and between places. If $R_t$ follows a stationary first-order autoregressive process with autocorrelation $\phi$, the environmental part of one-step prediction error falls to $(1-\phi^2)v_R$. If $J$ jurisdictions share environmental variance $v_R$ with mean pairwise correlation $\bar{c}$, the environmental variance of their average is given by:
\begin{equation}
  v_R\,\bar{c}+\frac{v_R(1-\bar{c})}{J},
  \label{eq:pooling}
\end{equation}
where pooling reduces this limit only to the extent that variation is not synchronised ($\bar{c}<1$).

\subsection{Reporting heterogeneity and the allocation of a surveillance budget}
\label{sec:bg-reporting}

Reporting completeness is rarely constant. Assumption H1 can be replaced by a reporting probability $p_t$ that varies between periods independently of $R_t$ and $n$, with mean $\rho$ and relative volatility $\tau=\mathrm{SD}(p_t)/\rho$, and with counts independent across periods conditional on $R_t$ and $p_t$ (H5).

\begin{proposition}[Reporting heterogeneity]
\label{prop:report-noise}
Under H2--H5, the error variance of $\Rhat_t=C_t/m$ about $\Rbar$ can be written as:
\begin{equation}
  V_H(m)=v_R+v_H+\frac{A_H}{m},\qquad
  v_H=\tau^2(\Rbar^2+v_R),\qquad A_H=\Rbar+\frac{\rho(1+\tau^2)(\Rbar^2+v_R)}{k},
  \label{eq:vH-AH}
\end{equation}
where $v_H$ is the variance added by reporting heterogeneity. The limiting variance rises to $v_R+v_H$, the balance scale falls to $m_{\times,H}=A_H/(v_R+v_H)<\mx$, and $v_R$ and $v_H$ cannot be separated using the mean and variance of the counts alone, which identify only $(1+\tau^2)(\Rbar^2+v_R)$.
\end{proposition}

Proposition~\ref{prop:report-noise} is the identification limit for the empirical part of this study. An intercept estimated from counts alone combines transmission and reporting variation unless reporting is measured separately.

The balance scale also determines how a fixed budget is best divided. Suppose that $T$ periods are observed with case volume $m$ in each, that each period costs $c_0$ and each reported case costs $c_1$, and that the budget is $\mathcal{B}=T(c_0+c_1m)$.

\begin{proposition}[Allocation of a fixed budget]
\label{prop:allocation}
Under H1--H4 with $R_t$ independent across periods, the variance of the mean of $\Rhat_t$ over $T$ periods is given by:
\begin{equation}
  \frac{v_R}{T}+\frac{A}{Tm},
  \label{eq:alloc-indep}
\end{equation}
which, under the budget constraint, is minimised at:
\begin{equation}
  m^*=\sqrt{\frac{c_0}{c_1}\,\mx},
  \label{eq:optimal-m}
\end{equation}
where $T^*=\mathcal{B}/(c_0+c_1m^*)$, subject to $T^*\ge1$. If $R_t$ follows a stationary first-order autoregressive process with autocorrelation $\phi$, the sampling term is unchanged and the environmental term becomes $v_RS(\phi,T)/T^2$ with
\begin{equation}
  S(\phi,T)=\frac{T(1+\phi)}{1-\phi}-\frac{2\phi(1-\phi^T)}{(1-\phi)^2}.
  \label{eq:ar-sum}
\end{equation}
\end{proposition}

The optimal case volume per period is the geometric mean of $c_0/c_1$ and the balance scale, the same form as the optimal number of secondary units in two-stage sampling \citep{cochran1977sampling}. Once $m$ exceeds $m^*$, additional budget reduces error more when spent on further periods than on further cases. With autocorrelated environments the optimum also depends on the budget, and Appendix~\ref{app:extensions} gives the threshold at which it switches.

Reporting heterogeneity can reverse the ranking of coverage and completeness as targets for investment.

\begin{corollary}[Coverage versus completeness]
\label{cor:switch}
If $\tau$ does not depend on $\rho$, increasing coverage $n$ reduces $V_H$ more than the same proportional increase in $\rho$. If reporting variation arises from binomial sampling of $\nu$ units with $\Var(p_t)=\rho(1-\rho)/\nu$, increasing $\rho$ is preferable whenever:
\begin{equation}
  m>\frac{\rho^2(\nu-1)}{k},
  \label{eq:r3-switch}
\end{equation}
where $\nu$ is the number of reporting units sampled per period (or $m>\rho^2\nu/k$ under a Beta prior $p_t\sim\mathrm{Beta}(\rho\nu,(1-\rho)\nu)$ with variance $\rho(1-\rho)/(\nu+1)$).
\end{corollary}
